\documentclass[11pt]{article}

\usepackage[review]{acl}
\usepackage{times}
\usepackage{latexsym}
\usepackage[T1]{fontenc}
\usepackage[utf8]{inputenc}
\usepackage{microtype}
\usepackage{inconsolata}

\title{SymTrust: Symbolic Methods for Trustworthy and Verifiable Reasoning}

\author{Anonymous Authors}
\definecolor{darkblue}{rgb}{0, 0, 0.5}
\hypersetup{colorlinks=true, citecolor=darkblue, linkcolor=darkblue, urlcolor=darkblue}
\newcommand{\pk}[1]{{\textcolor{red}{~(PK: #1)}}}
\newcommand{\tp}[1]{{\textcolor{orange}{~(TP: #1)}}}
\newcommand{\sr}[1]{{\textcolor{blue}{~(SR: #1)}}}

\begin{document}
\maketitle
\begin{abstract}
Large language models are now used in many important areas, making it essential to know when their outputs can be trusted \citep{chen-etal-2024-critical,zhang-etal-2024-comprehensive-survey,zeng-etal-2023-robotics}. However, correct answers may still follow flawed reasoning \citep{wang-etal-2025-chain}. Verifying how models reach their answers has therefore become necessary. Buidling on this, we ask a central question: how can we determine whether an answer is supported by reliable reasoning rather than merely appearing correct?
One promising direction is to use methods that make reasoning easier to inspect and verify. A growing body of work explores symbolic and structured methods, including rules, constraints, programs, graphs, knowledge bases, and diagrams. These methods can improve model reasoning and express relevant information and relations in forms that support evaluation and verification \citep{pan-etal-2023-logic,olausson-etal-2023-linc,pei-etal-2025-fover}. \textbf{SymTRUST} focuses on the connection between trustworthy AI and symbolic and structured methods. The workshop asks: how can these methods help us improve and verify model reasoning and build more reliable, understandable, and trustworthy AI systems? We invite researchers working across domains, modalities, and reasoning tasks to discuss new methods, evaluations, benchmarks, applications, and open challenges at the intersection of trustworthy AI, symbolic reasoning, and structured representations.
\end{abstract}

\section{Introduction}
Foundation models can now solve many difficult tasks involving language,
images, science, and decision-making. However, a correct answer alone does
not prove that a model followed a reliable reasoning process. Models can
reach correct answers despite mistakes in their intermediate reasoning
\citep{wang-etal-2025-chain}. They can also produce convincing explanations
that do not faithfully show what influenced their answers
\citep{turpin-etal-2023-language,shailya-etal-2025-lext}. Their performance
may change when a question is reworded, irrelevant information is added,
or the same problem is presented using another representation or modality
\citep{mirzadeh-etal-2024-gsm,dagan-etal-2025-cast,
rajpal-etal-2026-spatial}. Looking only at the final answer is therefore
not enough to decide whether a model's reasoning can be trusted. We need
methods that can verify whether an answer is supported by valid reasoning,
relevant evidence, and the requirements of the task.

Symbolic and structured methods offer a promising way to support this
verification. Symbolic methods represent knowledge and reasoning through
explicit concepts, relations, rules, constraints, or programs. Structured
methods organize information through graphs, tables, diagrams, grids,
schemas, or knowledge bases. Compared with free-form natural-language
explanations, these methods can make the information, relations, and rules
behind an answer more explicit. This makes it easier to inspect the
reasoning, identify failures, and check model outputs against external
tools, formal rules, supporting evidence, or known constraints
\citep{pan-etal-2023-logic,olausson-etal-2023-linc,
pei-etal-2025-fover}. In this way, symbolic and structured methods can help
us study and improve a central aspect of trustworthy AI: whether a model's
answer can be independently understood and verified.

TIES brings together research on trustworthy AI,
symbolic and neuro-symbolic reasoning, multimodal reasoning,
interpretability, formal verification, model probing, and human-centered
explanations. \sr{ do we bring trust first and then symbolic?: The workshop asks: \textbf{How can symbolic and structured
methods help us improve and verify model reasoning and build systems whose
behavior is more reliable, understandable, and trustworthy?}} We welcome
work from any domain, including law, finance, medicine, science,
mathematics, vision, robotics, and commonsense reasoning. Submissions may
present new methods, theories, datasets, benchmarks, evaluations,
applications, position papers, or analyses of model failures.
\paragraph{Workshop Goals and Distinction}

TIES focuses on a specific question at the intersection of trustworthy
AI and symbolic reasoning: when does making a model's reasoning explicit
actually make that reasoning more trustworthy? The workshop does not
assume that symbolic or structured representations are inherently
faithful. A generated proof may be logically valid but based on an
incorrect formalization; a graph may omit relevant evidence; and a
program may produce the correct answer while encoding the wrong
interpretation of the problem. TIES therefore studies both the promise
and the limitations of explicit representations as mechanisms for
improving, inspecting, and independently verifying model reasoning.

This focus distinguishes TIES from broader workshops on trustworthy NLP
and neuro-symbolic AI. Rather than addressing trustworthiness or
neuro-symbolic integration in general, TIES centers the relationship
between explicit structure and verifiable reasoning. It brings together
researchers developing symbolic methods with researchers studying
faithfulness, robustness, interpretability, evaluation, and
human-centered explanations.

\paragraph{Topics of Interest}

We invite archival and non-archival submissions on topics including, but
not limited to:

\begin{itemize}[leftmargin=12pt,noitemsep,topsep=1pt]
    \item generation and verification of formal or structured reasoning traces;
    \item neuro-symbolic and solver-augmented language models;
    \item rules, constraints, programs, graphs, knowledge bases, proofs, tables, and diagrams as reasoning representations;
    \item faithfulness and validity of structured explanations;
    \item autoformalization and failures in translating language into symbolic representations;
    \item reasoning across language, visual, spatial, and symbolic modalities;
    \item adaptive selection or switching between reasoning representations;
    \item formal verification and external tool use for model evaluation;
    \item robustness under paraphrases, counterfactuals, irrelevant information, and representation changes;
    \item evaluation metrics for reasoning validity, evidence support, calibration, and answer--trace consistency;
    \item human evaluation and human use of structured model explanations;
    \item multilingual and cross-domain structured reasoning;
    \item applications in law, medicine, finance, science, mathematics, robotics, and other high-stakes settings;
    \item analyses and negative results showing when explicit structure fails to improve reliability or trustworthiness.
\end{itemize}

\paragraph{Submission Types and Review}

TIES will welcome archival long papers of up to eight pages and short
papers of up to four pages, excluding references, following the ACL
formatting guidelines. We will also welcome two-page non-archival
extended abstracts describing ongoing work, position papers, recently
published work, datasets, system demonstrations, and negative results.
Non-archival submissions will be presented at the workshop but will not
be included in the archival proceedings.

We plan to accept both papers carrying completed ACL Rolling Review
reviews and papers submitted directly through OpenReview. Each direct
archival submission will receive three reviews. All accepted
contributions will be presented during one of two poster sessions.
Based on review quality, relevance, and topical diversity, selected
submissions will also be invited for oral presentation.

% ============================================================
% COMMENTED-OUT SHARED EVALUATION CHALLENGE
% Keep this section in the source in case it is added later.
% ============================================================

% \iffalse

% \paragraph{TIES Shared Evaluation Challenge}

% We plan to organize a shared evaluation challenge on
% \textit{Verifiable Reasoning Across Representations}. The challenge will
% evaluate whether systems can produce not only correct answers but also
% structured traces that provide independently checkable support for
% those answers.

% Each instance will contain a problem, a question, and identifiable
% supporting information. Equivalent instances may be presented in
% natural language and in a structured representation such as a set of
% rules, a graph, a table, or a spatial configuration. Participating
% systems will be asked to produce: (i) a final answer; (ii) a
% machine-readable reasoning trace; (iii) references to the evidence or
% constraints supporting individual steps; and, where appropriate, (iv)
% an uncertainty or abstention decision.

% The initial challenge will focus on spatial and rule-governed reasoning,
% where important parts of a trace can be verified programmatically.
% Evaluation will consider final-answer accuracy, structural validity,
% logical or executable correctness, evidence support, answer--trace
% consistency, and consistency across equivalent representations. Human
% annotations will provide reference answers, evidence links, and valid
% inference relations. Deterministic verification will be used wherever
% possible, with human evaluation used for semantic judgments that cannot
% be reduced to rule execution. LLM-based judging, if studied, will be
% reported only as a secondary evaluation rather than as the primary
% source of correctness.

% We anticipate participation from approximately 10--15 teams. Challenge
% results, baseline systems, annotation guidelines, and lessons about
% evaluating multiple valid reasoning traces will be discussed in a
% dedicated workshop session.

% \fi

\paragraph{Workshop Format}

TIES is proposed as a one-day workshop consisting of three invited
talks, selected oral presentations, two poster sessions, and a panel
discussion. The invited talks will provide complementary perspectives
on neuro-symbolic reasoning, trustworthy and interpretable AI, and the
evaluation of model reasoning. Oral presentations will highlight
selected archival and non-archival submissions.

The two poster sessions will provide substantial time for interaction
between participants and will ensure that every accepted submission
receives visibility. The concluding panel will bring together invited
speakers, organizers, and selected participants to address the
question: \textit{When does explicit structure provide genuine evidence
for trustworthy reasoning, and when does it merely create the
appearance of reliability?}

\paragraph{Tentative Invited Speakers}

We plan to invite three speakers whose expertise covers the central
perspectives represented by TIES:

\begin{itemize}[leftmargin=12pt,noitemsep,topsep=1pt]
    \item Hima Lakkaraju, Harvard University --- trustworthy and
          interpretable AI, explanation faithfulness, and reliable
          foundation models;
    \item Andr\'e Freitas, University of Manchester and Idiap Research
          Institute --- neuro-symbolic NLP, formal representations, and
          controlled inference;
    \item Isabelle Augenstein, University of Copenhagen --- trustworthy
          NLP, evidence-grounded evaluation, factuality, and
          explainability.
\end{itemize}

All listed speakers are tentative. The complimentary workshop
registration provided by ACL will be allocated to one invited speaker.
Other speakers may participate remotely or through independently
secured support; no additional travel funding is currently committed.

\paragraph{Expected Participation and Outreach}

As a first-edition workshop connecting several active research
communities, we expect approximately 60--80 participants at peak
attendance, including authors, invited speakers, and researchers from
neighboring areas. We anticipate approximately 20--25 archival
submissions and 10--15 non-archival extended abstracts.

The Call for Papers will be distributed through ACL mailing lists,
relevant Special Interest Groups, neuro-symbolic AI and trustworthy AI
communities, academic and industry research networks, and social-media
channels. Organizers will also contact groups working on formal
reasoning, model explanations, multimodal reasoning, verification, and
high-stakes applications.

\paragraph{Hybrid Participation Plan}

To support productive participation regardless of attendance mode,
invited and oral presentations will support live remote presentation
when required, with prerecorded talks collected as technical backups.
Remote participants will be able to submit questions through a
moderated channel and join the panel discussion.

Posters and supplementary materials will be made available through the
workshop website, with dedicated virtual discussion spaces for remote
attendees. Session chairs will monitor both in-person and online
questions to ensure that remote participation is not treated as a
separate or secondary track.

\paragraph{Diversity and Inclusion}

TIES will pursue diversity across gender, geography, ethnicity,
institution type, career stage, research methodology, and application
domain. We will seek balanced representation among organizers, invited
speakers, panelists, and Program Committee members, including
researchers who have not repeatedly served as keynote speakers or
panelists at recent major conferences.

Our outreach will include relevant ACL affinity groups, student
communities, researchers from underrepresented regions, and groups
working on multilingual and low-resource reasoning. The workshop will
welcome non-archival work and extended abstracts to reduce barriers for
early-stage researchers and participants whose work does not fit
standard archival formats. Early-career researchers will be considered
for oral presentations, panel participation, and poster spotlights.
Speakers will be encouraged to use accessible slides, and online
materials will be provided in formats compatible with assistive
technologies.

The workshop will follow the ACL anti-harassment policy. Organizers will
identify clear points of contact for reporting concerns and will act
promptly to maintain a respectful and inclusive environment.

\paragraph{Technical and Special Requirements}

TIES requires standard audiovisual equipment for presentations, poster
space, reliable internet access, and video-conferencing support for
hybrid participation. A room layout supporting poster sessions and
audience discussion would be beneficial. The workshop does not require
specialized on-site computing infrastructure.

\paragraph{Previous Editions}

TIES is a first-edition workshop. Although the workshop has not
previously been held under this title, the organizing team has research
and organizational experience spanning trustworthy AI,
neuro-symbolic learning, structured and multimodal reasoning,
explanation evaluation, and major NLP and AI venues. Relevant
individual experience is described in the supplementary organizer
information.

\paragraph{Tentative Schedule}

The proposed schedule for the one-day workshop is shown in
Table~\ref{tab:ties-schedule}. The schedule allocates substantial time
to poster presentations and discussion while balancing invited and
contributed research.

\begin{table}[t]
    \small
    \centering
    \renewcommand{\arraystretch}{1.08}
    \begin{tabular}{p{0.16\linewidth}p{0.72\linewidth}}
        \toprule
        Time & Session \\
        \midrule
        09:00--09:15 & Opening Remarks and Workshop Framing \\
        09:15--10:00 & Invited Talk 1 \\
        10:00--11:00 & Oral Presentations: Methods and Representations \\
        11:00--12:00 & Poster Session 1 and Coffee Break \\
        12:00--12:45 & Invited Talk 2 \\
        12:45--14:00 & Lunch Break \\
        14:00--15:00 & Oral Presentations: Evaluation and Applications \\
        15:00--16:00 & Poster Session 2 and Coffee Break \\
        16:00--16:45 & Invited Talk 3 \\
        16:45--17:45 & Panel Discussion and Audience Q\&A \\
        17:45--18:00 & Closing Remarks \\
        \bottomrule
    \end{tabular}
    \caption{Tentative schedule for the one-day TIES workshop.}
    \label{tab:ties-schedule}
\end{table}

% ============================================================
% SUPPLEMENTARY INFORMATION
% The material below belongs in the two supplementary pages.
% ============================================================

\newpage
\section*{Supplementary Information}

\paragraph{Organizing Committee}

\subparagraph{Shreya Rajpal}
\phantom{ }\\
\indent Affiliation: Michigan State University\\
\indent E-mail: \texttt{[EMAIL]}\\
\indent Website: \url{[WEBSITE]}

Shreya Rajpal is a Ph.D. student in Computer Science at Michigan State
University. Her research focuses on trustworthy AI, explanation
evaluation, neuro-symbolic reasoning, and reasoning across language and
structured or visual representations. 
% Her work includes the LExT
% framework for evaluating the trustworthiness of natural-language
% explanations and research on adaptive switching between linguistic and
% symbolic representations for spatial reasoning. 
She has held technical
leadership roles in student and professional research communities and
has experience presenting work at ACL-affiliated and interdisciplinary
AI venues.

\subparagraph{Parisa Kordjamshidi}
\phantom{ }\\
\indent Affiliation: Michigan State University\\
\indent E-mail: \texttt{kordjams@msu.edu}\\
\indent Website:
\url{https://www.cse.msu.edu/~kordjams/}

Parisa Kordjamshidi is an Associate Professor of Computer Science and
Engineering at Michigan State University. Her research focuses on
natural language processing, multimodal reasoning, and neuro-symbolic
methods for integrating structured knowledge with neural models. She
has organized workshops at ACL, EMNLP, NAACL, AAAI, IJCAI, and NeurIPS
and has held senior organizational and editorial roles in major NLP and
AI venues. Her expertise provides the workshop with both a strong
technical foundation and extensive experience in community
organization.

% Add at least two external organizers. Ideally include:
% 1. a senior researcher in trustworthy or interpretable AI;
% 2. a researcher in formal or neuro-symbolic reasoning;
% 3. optionally, an industry or high-stakes-domain researcher.

% \subparagraph{[External Organizer 1]}
% \phantom{ }\\
% \indent Affiliation: [AFFILIATION]\\
% \indent E-mail: \texttt{[EMAIL]}\\
% \indent Website: \url{[WEBSITE]}

% [Two to five sentences describing the organizer's research expertise,
% connection to TIES, and experience organizing workshops or related
% events.]

% \subparagraph{[External Organizer 2]}
% \phantom{ }\\
% \indent Affiliation: [AFFILIATION]\\
% \indent E-mail: \texttt{[EMAIL]}\\
% \indent Website: \url{[WEBSITE]}

% [Two to five sentences describing the organizer's research expertise,
% connection to TIES, and experience organizing workshops or related
% events.]

\paragraph{Program Committee and Review Capacity}

We expect approximately 25--35 archival submissions. To provide three
reviews per paper while assigning no reviewer more than three papers,
we will recruit approximately 35--40 Program Committee members. The
committee will represent expertise in formal reasoning, neuro-symbolic
learning, interpretability, multimodal reasoning, evaluation, and
high-stakes applications. It will also reflect geographic,
institutional, demographic, and career-stage diversity.

% \begin{itemize}[leftmargin=12pt,noitemsep,topsep=1pt]
%     \item [NAME, AFFILIATION] --- [agreed/tentative]
%     \item [NAME, AFFILIATION] --- [agreed/tentative]
%     \item [NAME, AFFILIATION] --- [agreed/tentative]
%     \item [NAME, AFFILIATION] --- [agreed/tentative]
%     \item [NAME, AFFILIATION] --- [agreed/tentative]
%     \item [Continue until approximately 35--40 members]
% \end{itemize}

\paragraph{Review Platform}

TIES will use OpenReview for direct workshop submissions and will also
consider papers submitted with completed ACL Rolling Review reviews.
Direct archival submissions will receive three reviews. Non-archival
extended abstracts will receive a lighter review focused on relevance,
clarity, and potential to contribute to workshop discussion. Conflicts
of interest will be managed according to ACL policies.

% \paragraph{Concurrent Proposals}

% % Replace this statement after checking with every organizer.
% No organizer is currently involved in another workshop proposal for the
% 2027 ACL conference cycle.



\bibliography{references}
\bibliographystyle{acl_natbib}

\end{document}